\documentclass{njupre/njupre}
\title[组会报告]{ 组会报告 }
\author[朱子航]{\texorpdfstring{朱子航 \\ \smallskip \textit{522022150087@smail.nju.edu.cn}}{}}
\date[\today]{\texorpdfstring{2023-03-20}{}}
\begin{document}
\begin{frame}
    \titlepage
\end{frame}
\begin{frame}
    \frametitle{目录}
    \tableofcontents
\end{frame}

\begin{frame}
    \frametitle{Gaussian Splatting = EM + Optim}
    正向：渲染过程
    \begin{itemize}
        \item 从3D高斯场景描述转为相机平面上的2D高斯分布，本质是一个投影变换$y=Mx$
        \item 从2D高斯分布分块光栅叠合成图片过程的本质是一个加权平均$pix=\sum\limits_{i=1}\limits^{N} \alpha_i w_in_i$
    \end{itemize}  
    逆向：估计过程
    \begin{itemize}
        \item 加权平均$pix=\sum\limits_{i=1}\limits^{N} \alpha_i w_in_i$的反向：从目标像素样本估计2D高斯混合模型分布（EM）
        \item $y=Mx$的反向：从某个未知3D高斯的若干2D投影样本估计原本的3D高斯分布，本质是给定若干$M_1,M_2,\dots,M_k$和$y_1,y_2,\dots,y_k$，可以数学化为求解优化问题$\arg\min_X \sum\limits_{i=1}\limits^{k} (M_iX-y_i)^2$（最小二乘法）
    \end{itemize}   
\end{frame}


\begin{frame}[allowframebreaks]{Reference}
    \bibliography{ref}
    \bibliographystyle{plain}
\end{frame}
\end{document}